\documentclass[10pt,a4paper]{amsart}
\usepackage[margin=19mm]{geometry}
\usepackage{amsmath,amssymb,mathtools}
\usepackage[scr=rsfs,cal=euler]{mathalfa}
\usepackage{xfrac}
\usepackage[unicode,hidelinks,bookmarks=false]{hyperref}
\newcommand{\ie}{i.e.\ }
\newcommand{\cf}{cf.\ }
\newcommand{\resp}{respectively\ }
\newcommand{\Subsection}{\subsection}
\providecommand{\nobreakdash}{\nobreak\hbox{-}}
\setlength{\emergencystretch}{2em}
\multlinegap0pt
\usepackage{tikz}
\usepackage{float}
\usepackage[shortlabels]{enumitem}
\usepackage{amssymb}
\usepackage{mathtools}
\newtheorem{lemma}{Lemma}
\title{A base-$8$ upper bound for planar peeling sequences}
\author{Andr\'e Hisatsuga, Griffin Johnston, Rafael Miyazaki}
\date{Source: arXiv:2609.13122v1 (CC BY 4.0)}
\begin{document}
\maketitle

\noindent\emph{Source and licence.} A base-$8$ upper bound for planar peeling sequences. Version arXiv:2609.13122v1 (CC BY 4.0), distributed by arXiv under the Creative Commons Attribution 4.0 International licence, http://creativecommons.org/licenses/by/4.0/, as recorded in the arXiv licence metadata for this exact version and shown on the abstract page https://arxiv.org/abs/2609.13122v1. Full licence notice: \texttt{licenses/arxiv-2609.13122-CC-BY-4.0.txt}.\par\smallskip

\noindent\textit{Editorial scope.} Counted unit: the article's lemma lem:weighted-kraft (source0.tex lines 570--581). The article's complete original proof of it is delivered with it (source0.tex lines 583--604, one \texttt{proof} environment of 22 lines). No mathematics is written, repaired or re-derived: the statement and the proof are the article's own lines, in the article's own notation, and no retained line is substituted or re-flowed.  No further source-local context is needed: every cross-reference of every retained line resolves inside the two delivered spans.  Nothing was omitted from any retained span, so the package carries no omission record.  No retained line contains a citation, so the wrapper carries no bibliography and no bibliography entry.  No retained line contains a figure environment, an image-inclusion command or drawing code, so no image is delivered and the package declares no asset dependency.  Editorial formatting only, in addition: an AMS \texttt{amsart} preamble that restates the article's own declarations and macros used by the retained lines, namely the environments \texttt{lemma} on separate counters, together with the standard packages those lines need.  The retained environments print in this document as Lemma 1 in order of appearance, whereas the article prints the same items with the numbering of its own full document.  The complete native source of the article as distributed in the arXiv e-print for this exact version is bundled unchanged as \texttt{sources/210183/source0.tex}.
\begin{lemma}\label{lem:weighted-kraft}
Let $T$ be a finite rooted tree. 
For every edge $v\to w$ of $T$, assign a
nonnegative weight $p(v,w)$, and suppose that for every vertex $v$,
\begin{equation*}
    \sum_{w:\, v\to w} p(v,w)\le 1.
\end{equation*}

For a vertex $u\in V(T)$, let its mass be the product of the weights
of the edges along the unique path from the root to $u$, with the root
having mass $1$. Then the set of leaves of $T$ has total mass at most $1$.
\end{lemma}

\begin{proof}
Let $\mu(u)$ denote the mass of a vertex $u$. Thus $\mu(r)=1$ for the
root $r$, and if $w$ is a child of $v$, then
$\mu(w)=\mu(v)p(v,w)$.
Consequently, for every vertex $v$,
\begin{equation*}
      \sum_{w:\,v\to w}\mu(w)
    =
    \mu(v)\sum_{w:\,v\to w}p(v,w)
    \le \mu(v).
\end{equation*}

Starting at the root and successively
replacing each non-leaf vertex $v$ by its children in
$T$ cannot increase the total mass, by the inequality
above. After finitely many such replacements, the resulting set of
vertices is precisely the set of leaves of $T$. Hence
\begin{equation*}
    \sum_{u \text{ leaf of }T}\mu(u)\le \mu(r)=1. \qedhere
\end{equation*}
    
\end{proof}

\end{document}
